% page 214 of the facsimile (image p-213 of the scan)
% block 162.6 x 233.3 mm ; left margin 8.0 mm
\pgc{-12.700mm}{0.000mm}{
\l{\marge[79.564mm]{23.347mm}{\ornamento{12.774mm}{ornaments/letters/letter-H.png}}\cel{2}206}
\l{}
\l{}
\l{}
\l{}
\l{}
\l{}
\l{}
\l{}
\l{ha !\cel{3}Interjeciono qua indikas ke onu quik expresos ula}
\l{\cel{3}emoco tre vivaca di psiko (joyo, doloro, iraco, e c.)}
\l{}
\l{\sou{hab}ila.\cel{2}(\sou{pri persono}). Apta sucesar en to quon lu agas o}
\l{\cel{3}facas. - DEFIS.}
\l{}
\l{habilitar. (\sou{trans.)}\cel{2}(\sou{yuro-cienco}). Igar kapabla legale pri}
\l{\cel{3}agor ulo. - DEFIS.}
\l{}
\l{\sou{habitar}.\cel{2}(\sou{trans.}) Rezidar permanante en domo propra o fixa.}
\l{\cel{3}- EFIS.}
\l{}
\l{\sou{habitaklo}. (\sou{navo)}. Armoro, sur la kastelo avana, avan la gu-}
\l{\cel{3}vernilo-roto, en qua enklozesas la busolo, la horlojo e la}
\l{\cel{3}lanternego. - F.}
\l{}
\l{\sou{hachar}. (\sou{trans.}) I. Tranchadar a peceti, per hakilo, hako-}
\l{\cel{3}kultelo, e c. ( (\sou{anke metaf.}) . II. Provizar (desegnuro, e c.}
\l{\cel{3}per mikra streki, paralela o qui interkrucumas. - EF.}
\l{}
\l{\sou{hagiografo}. Autoro qua naracas la vivo-maniero e la agi di}
\l{\cel{3}la santi (di la eklezio katolika). - DEFIS.}
\l{}
\l{\sou{hagiografio}. Cienco qua traktas la vivo-maniero e la agi di la}
\l{\cel{3}santi (di la eklezio katolika).}
\l{}
\l{\sou{hakar}. (\sou{trans., per)}\cel{1}. Abatar, tranchar, fendar, per instru-}
\l{\cel{3}mento fera, lamo larja ye formo di triangulo kurvolinea ;}
\l{\cel{3}konvexa, ye lua bazo (clqua esas tranchiva) ; konkava, ye}
\l{\cel{3}lua lateri, e qua esas fixigita an mancho. - DEFIS.}
\l{}
\l{\sou{halo}. I. Chambro, salono, extreme vasta. - II. Granda placo}
\l{\cel{3}tektoza, ube esas la merkato di la manjebli. - DEFS.}
\l{}
\l{\sou{halbardo}. Shaft-armo, kun mancho longa, ferlamo tranchiva e}
\l{\cel{3}pintoza, plus du ferlami laterala (ali), de qui un esas}
\l{\cel{3}krecentatra, la altra pinta. - DEFIRS.}
\l{}
\l{\sou{halofito}. (\sou{biol.}) Nomo generala di la vejetanti qui kreskas}
\l{\cel{3}naturale en sulo tre mar-saloza, sive an maro, sive apud}
\l{\cel{3}salini.}
\l{}
\l{halogena. Nomo donita da Berzelius ad irga ek la korpi ye qui}
\l{\cel{3}konsistas la familio di kloro : fluoro, kloro, bromo, iodo,}
\l{\cel{3}qui esas apta formacar sali per kombinar su kun la metali.}
}
